% answers_short.tex —— 2012 年·山东卷·文科 · 简短答案速查
% 由 tools/make_short_answers.py 自 answers.tex 的 \gksolhead/\gkanswer 提取，
% 并与 2012/2012_数学真题_山东文_答案要点.md 逐题比对后生成，勿手改（改源 MD 后重跑）。

\begin{gknotice}
\textbf{使用说明}：本节只给最终答案便于快速核对，完整推导见解析版。答案为编写者独立推导，\textbf{非官方标准答案}。
\end{gknotice}

\gkpart{选择题}
\begin{center}\begin{tabular}{ccccccccccccc}
\toprule
题号 & 1 & 2 & 3 & 4 & 5 & 6 & 7 & 8 & 9 & 10 & 11 & 12 \\
\midrule
答案 & A & C & B & D & C & A & B & A & B & D & D & B \\
\bottomrule
\end{tabular}\end{center}

\gkpart{填空题}
\begin{itemize}[leftmargin=2.6em]
  \item[\textbf{13.}] $\dfrac16$。
  \item[\textbf{14.}] $9$。
  \item[\textbf{15.}] $\dfrac14$。
  \item[\textbf{16.}] $(2-\sin2,\ 1-\cos2)$。
\end{itemize}

\gkpart{解答题}
\begin{itemize}[leftmargin=2.6em]
  \item[\textbf{17.}] (1) 证明见解析；(2) $\dfrac{\sqrt7}4$
  \item[\textbf{18.}] (1) $\dfrac3{10}$；(2) $\dfrac8{15}$
  \item[\textbf{19.}] (1) 证明见解析；(2) 证明见解析
  \item[\textbf{20.}] (1) $a_n=7n$；(2) $S_m=\dfrac{7(49^m-1)}{48}$
  \item[\textbf{21.}] (1) $\dfrac{x^2}4+y^2=1$；(2) 最大值为 $\dfrac{2\sqrt5}5$，此时 $m=0$ 或 $m=\pm\dfrac53$
  \item[\textbf{22.}] (1) $k=1$；(2) 递增区间 $(0,1)$，递减区间 $(1,+\infty)$；(3) 证明见解析
\end{itemize}
